\subsection{Delimitación}
% Tomar en cuenta los siguientes parámetros: 
% delimitación de contenidos (área académica, línea de investigación), 
% delimitación espacial (lugar donde se realizará la investigación), 
% delimitación temporal (período en el cual se desarrollará el Trabajo de Titulación).

Este trabajo de titulación se desarrolla dentro del área de la Ingeniería de Software y se enfoca en la línea de investigación de Sistemas de Información y Gestión de Datos. El proyecto se centra en estudiar, diseñar y construir un sistema de información centralizado para ordenar y unir el flujo de trabajo tanto de los doctores como del personal administrativo en un consultorio dental pequeño. Los temas específicos que marcan el límite técnico de la herramienta incluyen el análisis de las tareas diarias usando técnicas de diagramación de procesos, la recolección de lo que necesita el negocio y lo que exige la ley, y el diseño de una base de datos que junte las historias clínicas, el calendario de citas, las recetas médicas y los cobros en un solo lugar. Finalmente, se incluye la instalación del sistema y la medición de cuánto tiempo se ahorra comparando el trabajo antes y después de usar el software. Este proyecto no incluye conectar el sistema con plataformas externas del Ministerio de Salud Pública ni crear funciones de telemedicina o herramientas avanzadas para radiografías.

La investigación se realiza por completo en las instalaciones del centro odontológico Bellidental, un consultorio privado pequeño que está ubicado en la ciudad de Baños, en la provincia de Tungurahua, Ecuador. Todas las partes del proyecto, como hacer las entrevistas preliminares, observar cómo atienden a los pacientes, instalar el sistema y medir los tiempos de trabajo, se hacen de forma presencial dentro de este consultorio. Para el estudio se toma en cuenta al personal que trabaja todos los días ahí, el cual está conformado por un odontólogo general, un odontólogo especialista y una secretaria.

El tiempo que tomará desarrollar este trabajo de titulación va desde octubre de 2026 hasta enero de 2027. En estos meses se cumplirán todas las fases del proyecto. Primero se revisará cómo trabaja la clínica hoy en día y qué características debe tener el software. Después se diseñará el modelo de información con todas las reglas del consultorio. Por último, se pondrá a funcionar el sistema para comparar los tiempos que toma hacer el papeleo de los pacientes. Las mediciones de tiempo antes de tener el programa se tomarán al inicio del proyecto en Julio de 2026, y las pruebas finales para ver los resultados se harán al terminar el periodo en Noviembreero de 2026.